In this full-montage setup, per-electrode detection quality was concentrated at the frontopolar pair rather than distributed across the frontal region.
On Internal, Fp2 and Fp1 reached $F_1$ values of 86.82\% and 86.32\%, respectively, before performance fell at the neighbouring frontal sites AF4 and AF3 to 80.32\% and 73.76\%, with F3, and F4 at 45.65\%, and 45.23\%.
On Cao2018, Fp1 and Fp2 reached $F_1$ values of 79.68\% and 78.45\%, respectively, whereas the neighbouring frontal sites F3 and F4 were lower at 59.51\% and 57.99\%.

The central electrodes showed a dataset-specific asymmetry. On Internal, C3 reached $F_1=61.06\%$ while C4 reached $F_1=1.85\%$, whereas on Cao2018 C3 and C4 were close at $F_1=31.08\%$ and $F_1=30.18\%$, respectively, so the pronounced C3--C4 asymmetry was specific to Internal.

The posterior electrodes performed markedly worse than the frontal and central sites, and the parietal and occipital groups also differed from each other between datasets. On Internal, the four parietal electrodes CP1, CP2, CP5, and CP6 were uniformly near floor, with $F_1$ confined to a narrow 1.75\%--2.19\% range, and the eight occipital electrodes P7, P3, PO3, O1, PO4, O2, P4, and P8 showed only slightly more spread, from $F_1=1.77\%$ at P4 to $F_1=3.60\%$ at PO3. On Cao2018, the two parietal electrodes were more informative but still well below the frontal and central regions, with CP3 reaching $F_1=17.84\%$ and CP4 reaching $F_1=19.56\%$. The Cao2018 occipital group was more heterogeneous than its Internal counterpart. O1 and O2 sat near floor at $F_1=3.10\%$ and $3.33\%$, T5 and T6 were intermediate at $F_1=9.26\%$ and $8.30\%$, and P3 and P4, the two most anterior electrodes assigned to this region, reached $F_1=12.80\%$ and $12.67\%$, respectively, so the region's mean $F_1$ was driven up by P3 and P4 rather than by the electrodes located over true occipital cortex.
